%% ---- sezione A: in-breve+problema ----
\begin{riquadro}{La tesi in un paragrafo}
Il lavoro, svolto in collaborazione con Intesa Sanpaolo, affronta i limiti delle matrici di correlazione empiriche, affette da rumore statistico e vincolate all'unica traiettoria offerta dalla storia. A partire dai rendimenti giornalieri di 362 titoli dell'S\&P 500 con quotazioni ininterrotte nel periodo 2000--2020, la tesi segue un percorso incrementale, dal benchmark lineare con la PCA all'autoencoder lineare (AE lineare), all'autoencoder profondo (AE) e infine al variational autoencoder (VAE). Ogni modello è addestrato non sulla matrice grezza ma sul vettore del fattore di Cholesky, così che la ricomposizione $L L^{\top}$ garantisca per costruzione la semidefinita positività. L'obiettivo è generare scenari sintetici delle matrici di correlazione a un orizzonte nominale di 21 giorni lavorativi, costruiti ricampionando transizioni storiche dello spazio latente del VAE, per stimare la possibile variabilità delle correlazioni a quell'orizzonte in ottica di rischiosità. Il ricampionamento è un block-bootstrap degli incrementi latenti giornalieri, ancorato allo stato di mercato osservato all'ultima data del campione, in luogo del campionamento dal prior. Nella verifica aggregata tutti i 1000 scenari risultano conformi ai cinque fatti stilizzati delle matrici di correlazione finanziarie e definiti positivi (autovalore minimo fra 0.0101 e 0.0185), in linea con la semidefinita positività garantita per costruzione.
\end{riquadro}

\section{Il problema e l'obiettivo}

\paragraph{Correlazione e diversificazione.} Nella Modern Portfolio Theory di Markowitz (1952) il rischio di un portafoglio dipende soprattutto dal co-movimento dei titoli, catturato dalla matrice di correlazione: con correlazioni inferiori a 1 il rischio complessivo è strettamente inferiore alla somma dei rischi individuali (cap.~2, \S2.1). Basarsi sulla sola matrice storica significa però assumere interdipendenze immutate, mentre la struttura delle correlazioni è soggetta a continui cambiamenti, specie nei periodi di stress.

\paragraph{Tre limiti della matrice empirica.} Le criticità sono di tre ordini:
\begin{itemize}
\item la stima è rumorosa: l'accuratezza dello stimatore dipende dal parametro di spettro $q = N/T$ e, quando il numero di asset $N$ è confrontabile con la lunghezza $T$ della serie, la Random Matrix Theory (Bouchaud e Potters, 2011; Plerou et al., 2002) descrive la comparsa di autovalori artificialmente piccoli e grandi (\S2.2);
\item la storia offre una sola traiettoria e preclude configurazioni alternative o eventi estremi plausibili ma mai verificatisi nel campione (\S1.1);
\item i motori Monte Carlo classici del risk management usano un'unica matrice fissa lungo l'orizzonte: la dipendenza non può cambiare regime né rappresentare l'aumento brusco delle correlazioni nei ribassi (Longin e Solnik, 2001; Ang e Chen, 2002; \S2.3).
\end{itemize}

\paragraph{L'obiettivo.} Data la natura intrinsecamente stocastica dei mercati, la tesi giudica eccessivamente ambiziosa la previsione puntuale della matrice futura: l'obiettivo realistico è approssimare con insiemi di scenari la distribuzione delle configurazioni possibili. Il lavoro si propone quindi, come anticipato nel riquadro iniziale, di generare scenari sintetici delle matrici di correlazione a un orizzonte nominale di 21 giorni lavorativi (circa un mese), costruiti ricampionando transizioni storiche dello spazio latente, per stimare la possibile variabilità delle correlazioni a quell'orizzonte in ottica di rischiosità (\S2.1; \S8.1).

\paragraph{Perché i modelli generativi profondi.} Rispetto ai metodi tradizionali, che tendono a linearizzare le relazioni o a restare vincolati alla dinamica storica, i VAE (Kingma e Welling, 2013) mappano le matrici in uno spazio latente continuo da cui campionare stocasticamente. I ``pixel'' del titolo alludono alla filiazione di questi modelli dal dominio delle immagini: le matrici sono trattate come oggetti strutturati ad alta dimensionalità, pur con il fattore di Cholesky vettorizzato come input. I precedenti sono CorrGAN (Marti, 2020), che campiona matrici di correlazione con una GAN e ne verifica l'output sui cinque fatti stilizzati, e il lavoro di riferimento di Caprioli et al. (2024). Quest'ultimo addestra un VAE su matrici storiche e ne ricampiona le coordinate latenti, in analogia con la procedura del cap.~8 (\S1.2).

\paragraph{Il ruolo di Cholesky.} Ogni matrice di correlazione valida è per definizione semidefinita positiva (PSD, autovalori non negativi); la decomposizione di Cholesky richiede la condizione più forte di definita positività (rango pieno) e fornisce una matrice triangolare inferiore $L$ con $C = L L^{\top}$, usata per correlare vettori gaussiani $\mathbf{z} \sim \mathcal{N}(0, I_N)$ (\S2.3). Poiché l'output di un modello generativo non è garantito PSD, i modelli lavorano sul fattore $L$, la cui ricomposizione $L L^{\top}$ assicura la semidefinita positività per costruzione.

\paragraph{Obiettivi specifici e limiti dichiarati.} Gli obiettivi specifici (\S1.3) sono tre: una pipeline di pre-processing sull'S\&P 500 (2000--2020) con parametrizzazione di Cholesky; uno studio comparativo del potere ricostruttivo lungo il percorso PCA, AE lineare, AE profondo, con metriche di errore e topologiche al variare di $K$; l'integrazione probabilistica tramite VAE, campionando matrici inedite e verificandone la conformità ai cinque fatti stilizzati. Un limite è dichiarato fin dall'inizio: il filtro di continuità delle quotazioni introduce un survivorship bias (esclusione delle aziende fallite o uscite dall'indice), accettato per disporre di un panel bilanciato senza dati mancanti. Resta fuori perimetro la traduzione degli scenari in misure di rischio di portafoglio, indicata come primo sviluppo futuro (\S9.3).
